\section{Introduction}
\label{sec:intro}

The carbon intensity of electricity varies by hour and by location. Scheduling
computation to exploit that variation, by deferring work to cleaner hours or by
relocating it to cleaner grids, is now a well-established line of research, and
the reductions reported for it are large. They are also inconsistent. Studies
of temporal shifting report reductions that vary substantially with the region
examined \cite{wiesner2021letswait}, and studies of combined spatial and temporal
shifting report ranges wider still. Zanotto et
al.~\cite{zanotto2025usercentric}, whose formulation this paper takes as its
starting point, report a 79.25\,\% reduction against a carbon-agnostic baseline
in their most permissive configuration and 13.46\,\% in their most restrictive
one. Work examining the limits of these techniques has established that headline
figures depend heavily on the assumptions behind them
\cite{sukprasert2024limitations}, but the parameters responsible have not been
isolated individually.

The two configurations that produce the 13\,\% and the 79\,\% differ in two
respects at once: the permissive one allows fifty simultaneous jobs per region
where the restrictive one allows five, and it draws on a different set of
candidate regions. Neither parameter is varied alone, so the six-fold difference
cannot be attributed to either. The same study is explicit that the per-region
concurrency limit is consequential and that its estimation is left for future
work. This paper takes up that question, and finds that the answer accounts for
much of the range.

This paper argues that a substantial part of the variation is not a property of
the schedulers at all. It is produced by three choices in the evaluation which
are seldom reported, and each of which moves the reported figure by more than
the difference between the policies being compared.

The first is capacity. A scheduler that defers a request to a cleaner hour can
do so only if capacity is free in that hour. Evaluations that place a workload
into an effectively unconstrained pool measure the benefit of flexibility under
the most favourable condition available, and this condition is rarely stated.
We find that widening the deadline margin from 6\,h to 48\,h reduces emissions
by 32.1\,\% when the system runs at 13\,\% utilisation and by 0.9\,\% when it
runs at 82\,\%, a difference of more than thirty-fold produced entirely by how
heavily the system is loaded.

The second is the origin region assigned to a time-shifting policy. Such a
policy holds work in one region and moves it only in time, so its result is
partly a statement about that region. We find that the correlation between the
mean carbon intensity of the origin region and the emissions the policy achieves
is between $+0.941$ and $+0.998$, whereas the correlation with the region's
diurnal amplitude, the quantity the policy actually acts upon, is
indistinguishable from zero. Across the seven candidate origins in our primary
region set, the same policy yields results spanning a factor of 3.6. A single
reported figure is one draw from that spread, and the conventional choice of
origin selects the cleanest available grid.

The third is the composition of the candidate region set. Running an identical
workload and an identical scheduler at matched utilisation over two region sets,
we find that one reports approximately twice the reduction of the other at every
operating point. The set that reports more contains a region 4.7 times cleaner
than its nearest rival, so that a single destination is optimal for nearly every
request; the set that reports less does not.

These are not competing explanations for the spread in the literature but
complementary ones, and they share a mechanism. Across every configuration we
examine, the correlation between the fraction of the workload placed in the
busiest region and total emissions is below $-0.93$. Carbon-aware placement
policies work by concentrating load into the cleanest reachable grid; capacity
governs how much can be concentrated, the origin region governs where a
time-shifting policy starts from, and the region set governs how much is to be
gained by concentrating at all. The reductions these policies report are, to
first approximation, measurements of the concentration they achieve.

The consequence for practice is that the largest reductions are obtained in
precisely the configurations least likely to be deployable. At the lowest
utilisation we examine, the policies achieving the largest reductions place
between 95\,\% and 98\,\% of the workload in a single region.

Our contributions are as follows.

\begin{itemize}
\item We isolate the interaction between capacity utilisation and deadline
      margin, and show that the value of deadline flexibility is conditional on
      load rather than being a property of the workload or the scheduler
      (Section~\ref{sec:res-capacity}).
\item We show that the reported performance of a time-shifting policy is
      determined by its origin region rather than by the temporal structure it
      exploits, and quantify the resulting spread
      (Section~\ref{sec:res-home}).
\item We show that reported savings scale with the degree to which the candidate
      region set contains a dominant clean region, holding workload, scheduler
      and hardware fixed (Section~\ref{sec:res-regionset}).
\item We establish that these results hold across five workload samples, two
      seasons, two workload classes, two region sets and two host power models,
      and we report the accounting basis on which each is measured
      (Section~\ref{sec:res-robustness}).
\end{itemize}

We adopt the host platform, the simulator, the carbon-agnostic baseline, the
per-region concurrency constraint and the deadline-margin levels of Zanotto et
al.~\cite{zanotto2025usercentric}, so that differences between their reported
figures and ours cannot be attributed to any of those. Where we depart from them
is in varying one parameter at a time, in sweeping the origin region of the
time-shifting policy rather than fixing it, and in reporting emissions on an
explicitly stated accounting basis. The evaluation
uses published hourly emission factors \cite{maji2022carboncast} and a
production virtual machine trace \cite{cortez2017resourcecentral}, and each
experimental cell is checked to confirm that every virtual machine executed in
the region its planner selected.
